\chapter{Feynman-Diagramme und Quarkliniendiagramme}

\section{Ziel dieses Kapitels}

Dieses Kapitel behandelt die Darstellung von Teilchenreaktionen und Zerfaellen mit
Feynman-Diagrammen und Quarkliniendiagrammen.

Feynman-Diagramme sind keine Bilder von echten Teilchenbahnen.
Sie sind eine symbolische Darstellung von Wechselwirkungen, Anfangszustaenden,
Endzustaenden und Austauschteilchen.

\begin{notebox}
Dieses Kapitel behandelt die Stichwoerter:
\begin{itemize}
    \item Feynman-Diagramme,
    \item Teilchenlinien,
    \item Antiteilchenlinien,
    \item Vertices,
    \item Austauschteilchen,
    \item Photon,
    \item Gluon,
    \item $W^\pm$-Boson,
    \item $Z^0$-Boson,
    \item elektromagnetische Wechselwirkung,
    \item starke Wechselwirkung,
    \item schwache Wechselwirkung,
    \item Quarkliniendiagramme,
    \item Beta-Zerfall,
    \item Paarvernichtung,
    \item Teilchenerzeugung,
    \item Erhaltungssaetze an Vertices.
\end{itemize}
\end{notebox}

\section{Was ist ein Feynman-Diagramm?}

\begin{definitionbox}
Ein \emph{Feynman-Diagramm} ist eine graphische Darstellung eines Terms in der
Stoerungsrechnung einer Quantenfeldtheorie.

Es zeigt, welche Teilchen einlaufen, welche Teilchen auslaufen und welche
Austauschteilchen die Wechselwirkung vermitteln.
\end{definitionbox}

\begin{conceptbox}
Ein Feynman-Diagramm ist kein klassisches Bahn-Bild.
Die Linien zeigen nicht, dass ein Teilchen exakt so durch den Raum geflogen ist.

Das Diagramm ist eine Rechen- und Ordnungshilfe fuer einen quantenmechanischen
Prozess.
\end{conceptbox}

\begin{examtipbox}
In der Pruefung ist meist nicht die mathematische Auswertung des Diagramms wichtig,
sondern:
\begin{itemize}
    \item Welche Teilchen sind Anfangs- und Endzustand?
    \item Welches Austauschteilchen ist beteiligt?
    \item Welche Wechselwirkung ist beteiligt?
    \item Sind Ladung, Baryonenzahl und Leptonenzahl erhalten?
    \item Aendert sich ein Quark-Flavor?
\end{itemize}
\end{examtipbox}

\section{Achsen und Leserichtung}

Es gibt verschiedene Konventionen fuer Feynman-Diagramme.
In diesem Skript verwenden wir die einfache Konvention:

\begin{notebox}
Die Zeit laeuft von links nach rechts.

Einlaufende Teilchen stehen links.
Auslaufende Teilchen stehen rechts.
\end{notebox}

\begin{conceptbox}
Manchmal wird in Buechern oder Vorlesungsfolien die Zeitachse nach oben gezeichnet.
Das ist ebenfalls erlaubt.

Wichtig ist nicht die Orientierung des Blattes, sondern die korrekte Verbindung der
Teilchenlinien und die Erhaltungssaetze an jedem Vertex.
\end{conceptbox}

\section{Elemente eines Feynman-Diagramms}

\begin{definitionbox}
Ein \emph{Vertex} ist ein Wechselwirkungspunkt in einem Feynman-Diagramm.
An einem Vertex treffen Teilchenlinien und Austauschteilchenlinien zusammen.
\end{definitionbox}

\begin{conceptbox}
Typische Linien:
\begin{itemize}
    \item Fermionen werden als gerade Linien gezeichnet.
    \item Photonen werden oft als Wellenlinien gezeichnet.
    \item Gluonen werden oft als Spirallinien gezeichnet.
    \item $W$- und $Z$-Bosonen werden als spezielle Bosonlinien gezeichnet.
\end{itemize}
\end{conceptbox}

\begin{notebox}
In reinem LaTeX ohne spezielle Feynman-Pakete beschreiben wir Diagramme in diesem
Skript meist durch Reaktionsgleichungen und Vertex-Strukturen.

Das ist fuer das Verstaendnis und fuer viele Pruefungsfragen ausreichend.
\end{notebox}

\section{Austauschteilchen und Wechselwirkungen}

\begin{center}
\begin{tabular}{|p{4.0cm}|p{4.0cm}|p{5.2cm}|}
\hline
Wechselwirkung & Austauschteilchen & typische Wirkung \\
\hline
elektromagnetisch & $\gamma$ & wirkt auf elektrische Ladung \\
\hline
stark & $g$ & wirkt auf Farbladung von Quarks und Gluonen \\
\hline
schwach geladen & $W^+$, $W^-$ & kann Flavor und Ladung aendern \\
\hline
schwach neutral & $Z^0$ & schwache Wechselwirkung ohne Ladungsaenderung \\
\hline
\end{tabular}
\end{center}

\begin{examtipbox}
Zuordnung merken:
\[
    \gamma
    \quad\Rightarrow\quad
    \text{elektromagnetisch}.
\]
\[
    g
    \quad\Rightarrow\quad
    \text{stark}.
\]
\[
    W^\pm,\;Z^0
    \quad\Rightarrow\quad
    \text{schwach}.
\]
\end{examtipbox}

\section{Erhaltungssaetze an Vertices}

An jedem Vertex muessen die relevanten Erhaltungssaetze erfuellt sein.

\begin{notebox}
Immer erhalten sind:
\begin{itemize}
    \item elektrische Ladung,
    \item Energie und Impuls,
    \item Drehimpuls,
    \item Baryonenzahl,
    \item Leptonenzahl in den Standardprozessen dieser Vorlesung.
\end{itemize}
\end{notebox}

\begin{conceptbox}
Flavor-Quantenzahlen wie Strangeness sind nicht bei allen Wechselwirkungen erhalten.

Starke und elektromagnetische Wechselwirkung aendern den Quark-Flavor nicht.

Die schwache Wechselwirkung kann den Quark-Flavor aendern.
\end{conceptbox}

\begin{examtipbox}
Wenn ein Diagramm einen Flavorwechsel zeigt, zum Beispiel
\[
    d \rightarrow u
\]
oder
\[
    s \rightarrow u,
\]
dann ist die schwache Wechselwirkung beteiligt.
Meist tritt dann ein $W$-Boson auf.
\end{examtipbox}

\section{Elektromagnetischer Vertex}

Der elektromagnetische Vertex koppelt ein geladenes Fermion an ein Photon.

\begin{formulabox}
Typische Vertex-Struktur:
\[
    f \rightarrow f + \gamma.
\]
Dabei ist $f$ ein geladenes Fermion.
\end{formulabox}

\begin{conceptbox}
Das Photon koppelt an elektrische Ladung.
Ein geladenes Teilchen kann daher ein Photon emittieren oder absorbieren.

Das Photon aendert den Flavor des Fermions nicht.
\end{conceptbox}

\begin{examplebox}
Ein Elektron kann ein Photon emittieren:
\[
    e^- \rightarrow e^- + \gamma.
\]

Ladung:
\[
    -1 = -1 + 0.
\]

Leptonenzahl:
\[
    1 = 1 + 0.
\]

Die Erhaltungssaetze passen.
\end{examplebox}

\begin{examtipbox}
Elektromagnetische Wechselwirkung:
\begin{itemize}
    \item Photon als Austauschteilchen,
    \item elektrische Ladung noetig,
    \item kein Quark-Flavorwechsel,
    \item kein Neutrino am einfachen elektromagnetischen Vertex.
\end{itemize}
\end{examtipbox}

\section{Elektron-Positron-Annihilation}

Ein wichtiges Beispiel ist die Paarvernichtung von Elektron und Positron.

\begin{formulabox}
\[
    e^- + e^+ \rightarrow \gamma + \gamma.
\]
\end{formulabox}

\begin{conceptbox}
Elektron und Positron koennen annihilieren.
Dabei entstehen Photonen.

Ein einzelnes Photon als Endzustand ist fuer freie Teilchen im Vakuum wegen
Energie- und Impulserhaltung nicht ausreichend.
Daher treten bei der einfachen Annihilation meist zwei Photonen auf.
\end{conceptbox}

\begin{examplebox}
Ladung:
\[
    Q(e^-)+Q(e^+)
    =
    -1+1
    =
    0.
\]

Endzustand:
\[
    Q(\gamma)+Q(\gamma)=0+0=0.
\]

Die elektrische Ladung ist erhalten.
\end{examplebox}

\begin{examtipbox}
Paarvernichtung:
\[
    \text{Teilchen}+\text{Antiteilchen}
    \rightarrow
    \text{Bosonen oder andere erlaubte Endzustaende}.
\]
Bei Elektron-Positron-Annihilation ist elektromagnetisch oft
\[
    e^-+e^+\rightarrow \gamma+\gamma
\]
der einfachste Prozess.
\end{examtipbox}

\section{Virtuelles Photon}

In Streuung und Teilchenerzeugung tritt haeufig ein virtuelles Photon auf.

\begin{definitionbox}
Ein \emph{virtuelles Teilchen} ist ein internes Teilchen in einem Feynman-Diagramm,
das nicht direkt als freies Endteilchen beobachtet wird.
\end{definitionbox}

\begin{formulabox}
Ein typischer Prozess ist:
\[
    e^- + e^+
    \rightarrow
    \gamma^*
    \rightarrow
    \mu^- + \mu^+.
\]
\end{formulabox}

\begin{conceptbox}
Das Symbol
\[
    \gamma^*
\]
bezeichnet ein virtuelles Photon.

Es vermittelt die elektromagnetische Wechselwirkung zwischen Anfangs- und Endzustand.
\end{conceptbox}

\begin{examtipbox}
Ein Sternchen am Teilchensymbol bedeutet oft:
\[
    \text{virtuell oder angeregt}.
\]
Bei
\[
    \gamma^*
\]
ist ein virtuelles Photon gemeint.
\end{examtipbox}

\section{Starker Vertex}

Der starke Vertex koppelt Quarks an Gluonen.

\begin{formulabox}
Typische Vertex-Struktur:
\[
    q \rightarrow q + g.
\]
\end{formulabox}

\begin{conceptbox}
Gluonen koppeln an Farbladung.
Quarks tragen Farbladung und koennen daher Gluonen emittieren oder absorbieren.

Der starke Vertex aendert den Quark-Flavor nicht.
Ein Up-Quark bleibt ein Up-Quark, ein Strange-Quark bleibt ein Strange-Quark.
\end{conceptbox}

\begin{examplebox}
Ein Quark kann ein Gluon emittieren:
\[
    u \rightarrow u + g.
\]

Der Flavor bleibt gleich:
\[
    u \rightarrow u.
\]

Die Farbladung kann sich aendern, aber die Gesamtfarbladung ist erhalten.
\end{examplebox}

\begin{examtipbox}
Starke Wechselwirkung:
\begin{itemize}
    \item Gluon als Austauschteilchen,
    \item wirkt auf Farbladung,
    \item Quark-Flavor bleibt erhalten,
    \item Hadronen sind beteiligt,
    \item Prozesse sind typischerweise sehr schnell.
\end{itemize}
\end{examtipbox}

\section{Gluon-Selbstwechselwirkung}

Ein wichtiger Unterschied zwischen QED und QCD ist:
Gluonen tragen selbst Farbladung.

\begin{conceptbox}
Photonen sind elektrisch neutral.
Daher gibt es in der einfachen QED keinen direkten Photon-Photon-Vertex.

Gluonen tragen Farbladung.
Daher koennen Gluonen direkt miteinander wechselwirken.
\end{conceptbox}

\begin{notebox}
In hoeheren Ordnungen der starken Wechselwirkung treten Gluon-Selbstwechselwirkungen
auf.

Das ist ein wichtiger Grund fuer die besonderen Eigenschaften der QCD:
\begin{itemize}
    \item Confinement,
    \item asymptotische Freiheit,
    \item Jetbildung.
\end{itemize}
\end{notebox}

\begin{examtipbox}
Merke:
Photonen tragen keine elektrische Ladung.
Gluonen tragen Farbladung.

Deshalb wechselwirken Gluonen selbst stark miteinander.
\end{examtipbox}

\section{Schwacher geladener Strom}

Die schwache Wechselwirkung mit $W^\pm$-Bosonen kann Teilchenflavors aendern.

\begin{formulabox}
Typische Quark-Vertices:
\[
    d \rightarrow u + W^-,
\]
\[
    u \rightarrow d + W^+.
\]
\end{formulabox}

\begin{formulabox}
Typische Lepton-Vertices:
\[
    e^- \rightarrow \nu_e + W^-,
\]
\[
    \nu_e \rightarrow e^- + W^+.
\]
\end{formulabox}

\begin{conceptbox}
Ein $W$-Boson traegt elektrische Ladung.
Deshalb kann sich die Ladung des Fermions am Vertex aendern.

Die Gesamtladung bleibt aber erhalten.
\end{conceptbox}

\begin{examplebox}
Beim Vertex
\[
    d \rightarrow u + W^-
\]
gilt fuer die Ladung:
\[
    -\frac{1}{3}
    =
    +\frac{2}{3}
    +
    (-1).
\]

Die elektrische Ladung ist erhalten.
\end{examplebox}

\begin{examtipbox}
Wenn sich die elektrische Ladung eines Fermions am Vertex aendert, ist ein geladenes
$W$-Boson beteiligt.
\end{examtipbox}

\section{Beta-Minus-Zerfall auf Quarkebene}

Der Beta-Minus-Zerfall ist ein zentrales Beispiel fuer ein schwaches
Feynman-Diagramm.

\begin{formulabox}
Auf Nukleonenebene:
\[
    n \rightarrow p + e^- + \overline{\nu}_e.
\]
\end{formulabox}

\begin{formulabox}
Auf Quarkebene:
\[
    d \rightarrow u + W^-,
\]
\[
    W^- \rightarrow e^- + \overline{\nu}_e.
\]
\end{formulabox}

\begin{conceptbox}
Das Neutron hat Quarkinhalt:
\[
    n=udd.
\]

Das Proton hat Quarkinhalt:
\[
    p=uud.
\]

Beim Beta-Minus-Zerfall wird also eines der Down-Quarks in ein Up-Quark umgewandelt:
\[
    d \rightarrow u.
\]
\end{conceptbox}

\begin{examplebox}
Gesamtprozess:
\[
    udd
    \rightarrow
    uud + e^- + \overline{\nu}_e.
\]

Baryonenzahl:
\[
    1 \rightarrow 1+0+0.
\]

Leptonenzahl:
\[
    0 \rightarrow 0+1-1.
\]

Ladung:
\[
    0 \rightarrow +1-1+0.
\]
\end{examplebox}

\begin{examtipbox}
Beta-Minus-Zerfall:
\[
    n \rightarrow p + e^- + \overline{\nu}_e.
\]

Quarkebene:
\[
    d \rightarrow u + W^-.
\]
Das $W^-$ zerfaellt in:
\[
    W^- \rightarrow e^-+\overline{\nu}_e.
\]
\end{examtipbox}

\section{Beta-Plus-Zerfall}

Beim Beta-Plus-Zerfall wird ein Proton in ein Neutron umgewandelt.

\begin{formulabox}
Auf Nukleonenebene:
\[
    p \rightarrow n + e^+ + \nu_e.
\]
\end{formulabox}

\begin{formulabox}
Auf Quarkebene:
\[
    u \rightarrow d + W^+,
\]
\[
    W^+ \rightarrow e^+ + \nu_e.
\]
\end{formulabox}

\begin{conceptbox}
Das Proton hat Quarkinhalt:
\[
    p=uud.
\]

Das Neutron hat Quarkinhalt:
\[
    n=udd.
\]

Beim Beta-Plus-Zerfall wird also eines der Up-Quarks in ein Down-Quark umgewandelt:
\[
    u \rightarrow d.
\]
\end{conceptbox}

\begin{examtipbox}
Beta-Plus-Zerfall:
\[
    p \rightarrow n + e^+ + \nu_e.
\]

Quarkebene:
\[
    u \rightarrow d + W^+.
\]
\end{examtipbox}

\section{Elektroneneinfang}

Elektroneneinfang ist eng mit dem Beta-Plus-Prozess verwandt.

\begin{formulabox}
Auf Nukleonenebene:
\[
    p + e^- \rightarrow n + \nu_e.
\]
\end{formulabox}

\begin{formulabox}
Auf Quarkebene:
\[
    u + e^- \rightarrow d + \nu_e.
\]
\end{formulabox}

\begin{conceptbox}
Man kann den Prozess als Austausch eines $W$-Bosons verstehen.

Ein Proton im Kern faengt ein Huellenelektron ein.
Dadurch wird ein Proton in ein Neutron umgewandelt.
\end{conceptbox}

\begin{examtipbox}
Elektroneneinfang:
\[
    Z \rightarrow Z-1,
    \qquad
    A \rightarrow A.
\]

Im Kern wird ein Proton zu einem Neutron.
\end{examtipbox}

\section{Schwacher neutraler Strom}

Neben dem geladenen Strom mit $W^\pm$ gibt es den neutralen schwachen Strom mit
dem $Z^0$-Boson.

\begin{formulabox}
Typische Vertex-Struktur:
\[
    f \rightarrow f + Z^0.
\]
\end{formulabox}

\begin{conceptbox}
Das $Z^0$-Boson ist elektrisch neutral.
Es kann schwache Wechselwirkungen vermitteln, ohne die elektrische Ladung des
Fermions zu aendern.

Der Flavor bleibt am einfachen neutralen Strom erhalten.
\end{conceptbox}

\begin{examplebox}
Ein Beispiel fuer einen neutralen schwachen Prozess ist:
\[
    \nu_\mu + e^- \rightarrow \nu_\mu + e^-.
\]

Dieser Prozess kann ueber Austausch eines $Z^0$-Bosons laufen.
\end{examplebox}

\section{Quarkliniendiagramme}

Quarkliniendiagramme zeigen, wie die Quarks in Hadronen waehrend einer Reaktion
umgeordnet oder umgewandelt werden.

\begin{definitionbox}
Ein \emph{Quarkliniendiagramm} ist eine Darstellung, in der die Quarklinien der
beteiligten Hadronen explizit verfolgt werden.
\end{definitionbox}

\begin{conceptbox}
Quarkliniendiagramme sind besonders nuetzlich fuer:
\begin{itemize}
    \item Hadronenreaktionen,
    \item Mesonzerfaelle,
    \item Baryonzerfaelle,
    \item Erkennen von Quark-Flavor-Aenderungen,
    \item Unterscheiden starker, elektromagnetischer und schwacher Prozesse.
\end{itemize}
\end{conceptbox}

\begin{examtipbox}
Bei Quarkliniendiagrammen immer fragen:
\begin{itemize}
    \item Welche Quarks sind im Anfangszustand?
    \item Welche Quarks sind im Endzustand?
    \item Werden Quark-Antiquark-Paare erzeugt?
    \item Aendert sich ein Flavor?
    \item Ist die Reaktion stark, elektromagnetisch oder schwach?
\end{itemize}
\end{examtipbox}

\section{Quarklinien bei starker Erzeugung}

Bei starken Prozessen bleibt der Quark-Flavor erhalten.
Neue Quark-Antiquark-Paare koennen erzeugt werden.

\begin{examplebox}
Betrachte die starke Erzeugung seltsamer Teilchen:
\[
    \pi^- + p \rightarrow K^0 + \Lambda^0.
\]

Quarkinhalte:
\[
    \pi^- = d\overline{u},
\]
\[
    p=uud,
\]
\[
    K^0=d\overline{s},
\]
\[
    \Lambda^0=uds.
\]
\end{examplebox}

\begin{conceptbox}
Im Endzustand gibt es ein $s$-Quark und ein $\overline{s}$-Antiquark.
Diese koennen stark als Paar erzeugt werden:
\[
    g \rightarrow s\overline{s}.
\]

Die gesamte Strangeness bleibt erhalten:
\[
    S_{\mathrm{links}}=0,
\]
\[
    S_{\mathrm{rechts}}=+1-1=0.
\]
\end{conceptbox}

\begin{examtipbox}
Starke Produktion von Strangeness erfolgt paarweise:
\[
    s\overline{s}.
\]

Ein einzelner Strangeness-Wechsel ist dagegen schwach.
\end{examtipbox}

\section{Quarklinien bei schwachen Zerfaellen}

Bei schwachen Zerfaellen kann ein Quark-Flavor geaendert werden.

\begin{examplebox}
Betrachte:
\[
    \Lambda^0 \rightarrow p+\pi^-.
\]

Quarkinhalte:
\[
    \Lambda^0=uds,
\]
\[
    p=uud,
\]
\[
    \pi^-=d\overline{u}.
\]
\end{examplebox}

\begin{conceptbox}
Die Strangeness aendert sich:
\[
    S(\Lambda^0)=-1,
\]
\[
    S(p)+S(\pi^-)=0+0=0.
\]

Also:
\[
    \Delta S=+1.
\]

Das ist ein Hinweis auf schwache Wechselwirkung.
Auf Quarkebene kann zum Beispiel ein strange-Quark schwach zerfallen:
\[
    s \rightarrow u + W^-.
\]
Danach:
\[
    W^- \rightarrow d+\overline{u}.
\]
\end{conceptbox}

\begin{examtipbox}
Wenn ein strange-Hadron in nicht-seltsame Hadronen zerfaellt, ist der Prozess oft
schwach.
Der Grund ist:
\[
    s \rightarrow u
\]
oder
\[
    s \rightarrow d
\]
ist ein Flavorwechsel.
\end{examtipbox}

\section{Innere Linien und virtuelle Teilchen}

Interne Linien in Feynman-Diagrammen stellen virtuelle Teilchen dar.

\begin{definitionbox}
Eine \emph{interne Linie} ist eine Linie, die nicht im Anfangs- oder Endzustand
endet, sondern nur innerhalb des Diagramms vorkommt.
\end{definitionbox}

\begin{conceptbox}
Interne Teilchen muessen nicht dieselbe Energie-Impuls-Beziehung erfuellen wie reale
freie Teilchen.

Sie sind nicht direkt beobachtbar.
Man beobachtet nur Anfangs- und Endteilchen.
\end{conceptbox}

\begin{examtipbox}
Nicht sagen:
Das virtuelle Teilchen wurde direkt gemessen.

Richtig ist:
Das virtuelle Teilchen ist Teil der theoretischen Beschreibung des Prozesses.
\end{examtipbox}

\section{Hoehere Ordnungen}

Feynman-Diagramme koennen unterschiedliche Ordnungen haben.
Hoehere Ordnungen enthalten mehr Vertices oder Schleifen.

\begin{definitionbox}
Ein \emph{Diagramm hoeherer Ordnung} enthaelt zusaetzliche Vertices, interne Linien
oder Schleifen.
\end{definitionbox}

\begin{conceptbox}
Beispiele fuer Effekte hoeherer Ordnung sind:
\begin{itemize}
    \item Vakuumpolarisation,
    \item Selbstenergie-Korrekturen,
    \item Vertexkorrekturen,
    \item Gluonabstrahlung,
    \item Schleifendiagramme.
\end{itemize}
\end{conceptbox}

\begin{notebox}
Hoehere Ordnungen sind wichtig fuer Praezisionsrechnungen.
Fuer qualitative Pruefungsfragen reicht oft das fuehrende Diagramm.
\end{notebox}

\begin{examtipbox}
Mehr Vertices bedeuten meistens eine hoehere Ordnung in der Kopplungskonstante.
Solche Prozesse sind oft weniger wahrscheinlich, koennen aber fuer Praezisionsphysik
wichtig sein.
\end{examtipbox}

\section{Kopplungskonstanten und Wahrscheinlichkeit}

Die Staerke eines Prozesses haengt von der Kopplung an den Vertices ab.

\begin{conceptbox}
Jeder Vertex traegt einen Faktor, der mit der Kopplungsstaerke der jeweiligen
Wechselwirkung zusammenhaengt.

Qualitativ:
\begin{itemize}
    \item starke Wechselwirkung: grosse Kopplung bei niedrigen Energien,
    \item elektromagnetische Wechselwirkung: Kopplung ueber elektrische Ladung,
    \item schwache Wechselwirkung: bei niedrigen Energien durch schwere $W$- und
    $Z$-Bosonen stark unterdrueckt.
\end{itemize}
\end{conceptbox}

\begin{examtipbox}
Ein Prozess kann erlaubt sein, aber trotzdem selten.
Erhaltungssaetze sagen, ob ein Prozess moeglich ist.
Kopplungen und Phasenraum bestimmen, wie wahrscheinlich er ist.
\end{examtipbox}

\section{Typische Diagramm-Erkennung}

\begin{center}
\begin{tabular}{|p{4.0cm}|p{5.0cm}|p{5.0cm}|}
\hline
Beobachtung & wahrscheinliche Wechselwirkung & Grund \\
\hline
Photon beteiligt & elektromagnetisch & Photon koppelt an elektrische Ladung \\
\hline
Gluon beteiligt & stark & Gluon koppelt an Farbladung \\
\hline
$W^\pm$ beteiligt & schwach geladen & Ladung und Flavor koennen sich aendern \\
\hline
$Z^0$ beteiligt & schwach neutral & schwach ohne Ladungsaenderung \\
\hline
Neutrino im Endzustand & schwach & Neutrinos wechselwirken schwach \\
\hline
Strangeness nicht erhalten & schwach & Flavorwechsel \\
\hline
\end{tabular}
\end{center}

\section{Beispiel: Myonzerfall}

Der Myonzerfall ist ein rein leptonischer schwacher Prozess.

\begin{formulabox}
\[
    \mu^- \rightarrow e^- + \overline{\nu}_e + \nu_\mu.
\]
\end{formulabox}

\begin{conceptbox}
Auf Vertex-Ebene:
\[
    \mu^- \rightarrow \nu_\mu + W^-,
\]
\[
    W^- \rightarrow e^- + \overline{\nu}_e.
\]
\end{conceptbox}

\begin{examplebox}
Leptonenzahlen:

Links:
\[
    L_\mu=1,
    \qquad
    L_e=0.
\]

Rechts:
\[
    \nu_\mu
    \quad\Rightarrow\quad
    L_\mu=1,
\]
\[
    e^-+\overline{\nu}_e
    \quad\Rightarrow\quad
    L_e=1-1=0.
\]

Die Familien-Leptonenzahlen sind in dieser einfachen Darstellung erhalten.
\end{examplebox}

\section{Beispiel: Pionzerfall}

Geladene Pionen zerfallen schwach.

\begin{formulabox}
\[
    \pi^- \rightarrow \mu^- + \overline{\nu}_\mu.
\]
\end{formulabox}

\begin{conceptbox}
Quarkinhalt:
\[
    \pi^- = d\overline{u}.
\]

Das Quark-Antiquark-Paar kann ueber ein virtuelles $W^-$ annihilieren:
\[
    d+\overline{u}\rightarrow W^-,
\]
\[
    W^- \rightarrow \mu^-+\overline{\nu}_\mu.
\]
\end{conceptbox}

\begin{examtipbox}
Hadronische Anfangszustaende koennen schwach in Leptonen zerfallen, wenn ein
$W$-Boson vermittelt.
Das ist ein typisches Zeichen schwacher Wechselwirkung.
\end{examtipbox}

\section{Beispiel: Elektron-Proton-Streuung}

Elektron-Proton-Streuung kann elektromagnetisch durch Photonenaustausch beschrieben
werden.

\begin{formulabox}
\[
    e^- + p \rightarrow e^- + p.
\]
\end{formulabox}

\begin{conceptbox}
Das Elektron tauscht ein virtuelles Photon mit dem Proton aus:
\[
    e^- \rightarrow e^- + \gamma^*,
\]
\[
    p+\gamma^* \rightarrow p.
\]

Bei groesserem Impulsuebertrag kann die innere Struktur des Protons sichtbar werden.
Dann beschreibt man die Wechselwirkung auf Quarkebene.
\end{conceptbox}

\begin{examtipbox}
Elektron-Proton-Streuung ist wichtig fuer den experimentellen Nachweis der
Nukleonensubstruktur.
Das virtuelle Photon sondiert die elektrische Struktur des Protons.
\end{examtipbox}

\section{Systematisches Vorgehen beim Zeichnen}

\begin{examtipbox}
Gehe beim Zeichnen eines Feynman-Diagramms so vor:

\begin{enumerate}
    \item Schreibe Anfangs- und Endteilchen auf.
    \item Bestimme die beteiligte Wechselwirkung.
    \item Waehle das passende Austauschteilchen:
    \[
        \gamma,\quad g,\quad W^\pm,\quad Z^0.
    \]
    \item Pruefe an jedem Vertex die elektrische Ladung.
    \item Pruefe Baryonenzahl und Leptonenzahl.
    \item Pruefe, ob Flavor erhalten oder geaendert wird.
    \item Zeichne interne Teilchen als virtuelle Teilchen.
\end{enumerate}
\end{examtipbox}

\begin{conceptbox}
Ein richtiges Diagramm muss nicht kuenstlerisch perfekt sein.
Es muss die richtige Physik zeigen:
\begin{itemize}
    \item richtige Teilchen,
    \item richtige Vertices,
    \item richtige Austauschteilchen,
    \item richtige Erhaltungssaetze.
\end{itemize}
\end{conceptbox}

\section{Mini-Zusammenfassung}

\begin{notebox}
Die wichtigsten Punkte dieses Kapitels sind:
\begin{itemize}
    \item Feynman-Diagramme sind symbolische Darstellungen von Wechselwirkungen.
    \item Vertices sind Wechselwirkungspunkte.
    \item Interne Linien stehen fuer virtuelle Teilchen.
    \item Das Photon vermittelt elektromagnetische Wechselwirkung.
    \item Gluonen vermitteln starke Wechselwirkung.
    \item $W^\pm$- und $Z^0$-Bosonen vermitteln schwache Wechselwirkung.
    \item Elektromagnetische und starke Vertices aendern den Quark-Flavor nicht.
    \item Geladene schwache Vertices mit $W^\pm$ koennen Flavor aendern.
    \item Beta-Minus-Zerfall auf Quarkebene:
    \[
        d \rightarrow u + W^-.
    \]
    \item Beta-Plus-Zerfall auf Quarkebene:
    \[
        u \rightarrow d + W^+.
    \]
    \item Quarkliniendiagramme verfolgen Quarkflavors in Hadronenreaktionen.
    \item Starke Strangeness-Erzeugung erfolgt paarweise als $s\overline{s}$.
    \item Einzelne Strangeness-Aenderungen sind typisch fuer schwache Prozesse.
\end{itemize}
\end{notebox}

\section{Pruefungscheck}

\begin{examtipbox}
Nach diesem Kapitel solltest du folgende Fragen beantworten koennen:
\begin{enumerate}
    \item Was ist ein Feynman-Diagramm?
    \item Warum ist ein Feynman-Diagramm kein klassisches Bahn-Bild?
    \item Was ist ein Vertex?
    \item Was bedeutet eine interne Linie?
    \item Was ist ein virtuelles Teilchen?
    \item Welches Austauschteilchen vermittelt die elektromagnetische Wechselwirkung?
    \item Welches Austauschteilchen vermittelt die starke Wechselwirkung?
    \item Welche Austauschteilchen vermitteln die schwache Wechselwirkung?
    \item Was prueft man an jedem Vertex?
    \item Warum aendert ein Photon keinen Quark-Flavor?
    \item Warum aendert ein Gluon keinen Quark-Flavor?
    \item Wie sieht der Beta-Minus-Zerfall auf Quarkebene aus?
    \item Wie sieht der Beta-Plus-Zerfall auf Quarkebene aus?
    \item Was passiert beim Elektroneneinfang auf Quarkebene?
    \item Was ist ein schwacher neutraler Strom?
    \item Was zeigt ein Quarkliniendiagramm?
    \item Warum ist starke Strangeness-Erzeugung paarweise?
    \item Warum sind einzelne Strangeness-Aenderungen schwach?
    \item Was bedeutet hoehere Ordnung in einem Feynman-Diagramm?
    \item Woran erkennt man elektromagnetische, starke und schwache Diagramme?
\end{enumerate}
\end{examtipbox}